% AUTHOR: Amlal El Mahrouss
% PURPOSE: DN001.10: Notes for Several Formulas.

\documentclass[11pt, a4paper]{article}
\usepackage{graphicx}
\usepackage{listings}
\usepackage{amsmath,amssymb,amsthm}
\usepackage{xcolor}
\usepackage{hyperref}
\usepackage{mathtools}
\usepackage[margin=0.5in,top=1in,bottom=1in]{geometry}

\title{Third Notes for Several Formulas.}
\author{Amlal El Mahrouss\\\texttt{amlal@nekernel.org}}
\date{February 21, 2026}

\begin{document}

\maketitle

\begin{center}
	\rule[0.01cm]{17cm}{0.01cm}
\end{center}

\begin{abstract}
\begin{center}
    We'll develop new formulas in this paper, where we will cover some new formulas and properties.
\end{center}
\end{abstract}

\section{DEFINITIONS}

We assume the following to be true, and $\operatorname{General}(t), N(u)$ to be assumed defined:

\section{DEFINITION OF A GRAIN-FUNCTION:}

We define $\operatorname{Grain}(\alpha, \phi, n)$ as:

\begin{equation}
    \operatorname{Grain}(\alpha, \phi, n) \coloneqq \frac{\partial \operatorname{General}_{Grain}(\alpha)}{\partial \alpha} = \sum_{k=N(\phi)}^{n} \operatorname{General}_{Grain}(|\alpha|)
\end{equation}
For $\phi \geq 1, \forall \alpha, \phi \in \mathbb{R}, \forall n, \quad \operatorname{Normal}(\phi) \in \mathbb{R}$

\section{DEFINITION OF A U-FUNCTION:}

Let an U-Function in the form of:

\begin{equation}
    \operatorname{General}_{Grain}(x) \in \mathbb{R}
\end{equation}
Be $\operatorname{General}_{Grain}(x) \geq x, \quad x \geq 1$.

\section{DEFINITION OF A N-FUNCTION:}

Let an N-Function in the form of:

\begin{equation}
    \operatorname{Normal}(x) \in \mathbb{R}
\end{equation}
Be $\operatorname{Normal}(x) > 0$.

\end{document}
